% !PRINTED 129
\weberpage{129}
\begin{center}
\S.\,32.

\gesperrt{Theorem of Frobenius.}
\end{center}

\gesperrt{Frobenius} has established a theorem which may be regarded as a counterpart to the fourth Sylow theorem, and which may be stated thus\footnote{\gesperrt{Frobenius}, Sitzungsberichte der Berl. Akademie, 4. Mai 1893.}:

V. \gesperrt{Every group whose degree is a product of nothing but distinct primes is metacyclic.}

We therefore assume that the degree $n$ of a group $P$ is \gesperrt{divisible by no square number}. If $t$ is the greatest prime divisor of $n$, then by the \gesperrt{Cauchy-Sylow} theorem (\S.\,29) there is a divisor $T$ of degree $t$ of $P$, which, because $t$ is a prime, is a cyclic group, and thus consists of the elements

\[
T = 1, \; a, \; a^2, \; \ldots \; a^{t-1}
\]

If it could now be shown that under the assumptions made $T$ is a normal divisor of $P$, then one could apply the same mode of reasoning to the group $P/T$, whose degree is $n:t$, and, if $t_1$ is the second-greatest prime dividing $n$, would conclude that there is a normal divisor $N$ of the group $P/T$ of degree $t_1$. But from this it would then follow (\S.\,6, 2.) that $P$ has a normal divisor $T_1$ of degree $t\,t_1$, of which $T$ itself is in turn a normal divisor. By considering the group $P/T_1$ in the same way, and so on, one obtains a composition series of $P$:

\[
P, \; T_k, \; T_{k-1}, \; \ldots \; T_1, \; T, \; 1,
\]

in which the index series consists of the prime factors of $n$ arranged in ascending order of magnitude, and $P$ proves to be metacyclic. Everything therefore comes down to showing that $T$ is a \gesperrt{normal divisor} of $P$.

This theorem is easier to prove if one regards it as a special case of a more general one than if one attacks it directly. Under the assumptions made about $P$ and $n$, this more general theorem reads:

$\alpha$) \gesperrt{If $n = \mu\nu$ is decomposed into two factors and every prime divisor of $\nu$ is greater than every prime}
%CONT>>
